Public key cryptography solves a problem that symmetric encryption alone handles poorly: secure key establishment between parties that do not already share a secret. This chapter examined two major constructions that achieve this goal through different mathematical foundations.

The first construction, \textbf{ElGamal}, is based on exponentiation in a cyclic group and derives its security from the hardness of Diffie--Hellman-style problems, especially the \textbf{Decisional Diffie--Hellman assumption}. Its use of fresh randomness during encryption is essential: that randomness is what prevents ciphertexts from revealing message equality and is what makes semantic security possible.

The second construction, \textbf{RSA}, is built from a trapdoor permutation on arithmetic modulo a composite integer. Its correctness follows from Euler's theorem, while its practical security is tied to the hardness of factoring the modulus. RSA shows how elegant number-theoretic structure can yield an efficient public key primitive with a clear forward direction and a hidden inverse.

At the same time, the chapter showed that \textbf{textbook RSA is insecure}. Deterministic encryption, vulnerability to small-message attacks, and the multiplicative structure of the raw function all prevent it from serving as a secure encryption scheme on its own. A trapdoor permutation is therefore \emph{not} the same thing as secure public key encryption.

To obtain practical security, RSA must be combined with \textbf{padding}. Padding introduces randomness and structured preprocessing before exponentiation, removing the most obvious weaknesses of the textbook scheme. The older \textbf{PKCS\#1 v1.5} encoding improved matters significantly, but history showed that it could still fail under adaptive chosen-ciphertext conditions. More modern designs such as \textbf{OAEP} were introduced to close that gap, offering provable IND-CCA2 security in the random oracle model.

The \textbf{Bleichenbacher attack} is the clearest real-world illustration of this principle, and its long-lived descendants (DROWN, ROBOT) show that the danger was not a one-time bug but a structural pitfall. The attack undermined deployed systems not by breaking the underlying mathematics, but by exploiting a single bit of information leaked during decryption. In that sense, the chapter's most important lesson is broader than ElGamal or RSA alone: \textbf{cryptographic security depends on both theory and implementation}. Strong hardness assumptions are necessary, but they are not sufficient unless the complete scheme---encoding, protocol, and implementation---is designed and deployed with equal care.
